\documentclass[12pt,a4paper]{ctexbook}
\usepackage{geometry}
\geometry{margin=2.2cm}
\usepackage{setspace}
\setstretch{1.35}
\usepackage{hyperref}
\title{全宋诗选·poet.song.1000}
\author{古典文献整理}
\date{}
\begin{document}
\maketitle
\tableofcontents
\newpage
\section{詠月}
\textbf{朱貞白}\par\vspace{0.5em}
當塗當塗見，蕪湖蕪湖見。\par
八月十五夜，一似沒柄扇。\par

\section{詠螃蟹}
\textbf{朱貞白}\par\vspace{0.5em}
蟬眼龜形脚似蛛，未嘗正面向人趨。\par
如今飣在盤筵上，得似江湖亂走無。\par

\section{句  其一}
\textbf{朱貞白}\par\vspace{0.5em}
倒排雙陸子，稀插碧牙籌。\par
既似柿牛妳，又如鈴馬兜。\par
鼓搥并㩧箭，直是有來由。\par

\section{句  其二}
\textbf{朱貞白}\par\vspace{0.5em}
一夜小園開似雪。\par

\section{句}
\textbf{邵拙}\par\vspace{0.5em}
萬國未得雨，孤雲猶在山。\par

\section{句}
\textbf{羅穎}\par\vspace{0.5em}
嫚侮羣豪誇大度，可憐容得辟陽侯。\par

\section{魰}
\textbf{劉吉}\par\vspace{0.5em}
八珍一箸千金價，往往精庖賤惠文。\par
莫道形模大剛拙，剖珠也解獻殷勤。\par

\section{句}
\textbf{劉吉}\par\vspace{0.5em}
一箭不中鵠，五湖歸釣魚。\par

\section{送季大夫牧舒州}
\textbf{湯悅}\par\vspace{0.5em}
却下烏臺建隼旟，侯封歸去襲龍舒。\par
嚴霜尚滿辭天闕，甘雨看隨入境車。\par

\section{奉和聖制送鄧王牧宣城}
\textbf{湯悅}\par\vspace{0.5em}
千里陵陽同陝服，鑿門胙土寄親賢。\par
曙烟已別黄金殿，晚照重登白玉筵。\par
江上浮光宜雨後，郡中遠岫列窗前。\par
天心待報期年政，留與工師播管弦。\par

\section{早春寄華下同志}
\textbf{湯悅}\par\vspace{0.5em}
正是花時節，思君寢復興。\par
市沽終不醉，春夢亦無憑。\par
嶽面懸清雨，河心走濁冰。\par
東門一條路，離恨正相仍。\par

\section{鼎臣學士侍郎以東館庭梅昔翰苑之毫末今復半枯向時同僚零落都盡素髮垂領茲唯二人感舊傷懷發於吟詠惠然好我不能無言輒次來韻攀和}
\textbf{湯悅}\par\vspace{0.5em}
憶見萌芽日，還憐合抱時。\par
舊歡如夢想，物態暗還移。\par
素豔今無幾，朱顔亦自衰。\par
樹將人共老，何暇更悲絲。\par

\section{再次前韻代梅答}
\textbf{湯悅}\par\vspace{0.5em}
託植經多稔，頃筐向盛時。\par
枝條雖已故，情分不曾移。\par
莫向階前老，還同鏡裏衰。\par
更應憐墮葉，殘吹挂蟲絲。\par

\section{鼎臣學士侍郎楚金舍人學士以再傷庭梅詩同垂寵和清絕感歎情致俱深因成四十字陳謝}
\textbf{湯悅}\par\vspace{0.5em}
人物同遷謝，重成念舊悲。\par
連華得瓊玖，合奏發塤篪。\par
餘枿雖無取，殘芳尚獲知。\par
問君何所似，珍重杜秋詩。\par

\section{謁聖林}
\textbf{楊文郁}\par\vspace{0.5em}
悠悠往古繼來今，天地無窮照孔林。\par
兩到金絲堂下拜，門生無負百年心。\par

\section{句}
\textbf{駱仲舒}\par\vspace{0.5em}
張鴻詩在楞伽峽，韓愈碑留燕喜亭。\par

\section{偈}
\textbf{釋省念}\par\vspace{0.5em}
今年六十七，老病隨緣且遣日。\par
今年記却來年事，來年記著今朝日。\par

\section{偈}
\textbf{釋省念}\par\vspace{0.5em}
白銀世界金色身，情與非情共一真。\par
明暗盡時俱不照，日輪午後示全身。\par

\section{句  其一}
\textbf{釋省念}\par\vspace{0.5em}
一言截斷千江口，萬仞峰前始得玄。\par

\section{句  其二}
\textbf{釋省念}\par\vspace{0.5em}
萬里神光都一照，誰人敢並日輪齊。\par

\section{贈成務崇}
\textbf{孟嘏}\par\vspace{0.5em}
同呼碧嶂前，已是十餘年。\par
話別非容易，相逢不偶然。\par
多爲詩酒役，早免利名牽。\par
幸有歸真路，何妨學上玄。\par

\section{句}
\textbf{孟嘏}\par\vspace{0.5em}
詩酒獨游寺，琴書多寄僧。\par

\section{贈天台逸人}
\textbf{廖融}\par\vspace{0.5em}
移檜托禪子，携家上赤城。\par
拂琴天籟寂，欹枕海濤生。\par
雪白寒峰晚，鳥歌春谷晴。\par
又聞求桂楫，載月十洲行。\par

\section{題寺古檜}
\textbf{廖融}\par\vspace{0.5em}
何人見植初，老對梵王居。\par
山鬼暗栖托，樵夫難破除。\par
聲高秋漢迥，影倒月潭虛。\par
盡日無僧繞，清風長有餘。\par

\section{夢仙謠}
\textbf{廖融}\par\vspace{0.5em}
琪木扶疏繫辟邪，麻姑夜宴紫皇家。\par
銀河旌節摇波影，珠閣笙簫吸月華。\par
翠鳳引游三島路，赤龍齊別五雲車。\par
星稀猶倚虹橋立，擬就張騫搭漢槎。\par

\section{退宮妓}
\textbf{廖融}\par\vspace{0.5em}
神仙風格本難儔，曾從前皇翠輦游。\par
紅躑躅繁金殿暖，碧芙蓉笑水宮秋。\par
寶車鈿剝陰塵覆，錦帳香消畫燭幽。\par
一旦色衰歸故里，月明猶夢按梁州。\par

\section{謝翁宏以詩百篇見示}
\textbf{廖融}\par\vspace{0.5em}
高奇一百篇，見造化工全。\par
積思游滄海，冥搜入洞天。\par
神珠迷罔象，瑞玉匪雕鑴。\par
休嘆不得力，離騷千古傳。\par

\section{書伍彬屋壁}
\textbf{廖融}\par\vspace{0.5em}
圓塘綠水平，魚躍紫蒓生。\par
要路貧無力，深村老退耕。\par
犢隨原草遠，蛙傍塹籬鳴。\par
撥棹茶川去，初逢穀雨晴。\par

\section{贈王正己}
\textbf{廖融}\par\vspace{0.5em}
吟高鄙俗流，傲逸訪巢由。\par
古寺尋僧飯，寒巖衣鹿裘。\par
園桃山鼠嚙，崖蜜獵人偷。\par
遂這箇清性，浮生不苟求。\par

\section{贈狄渙}
\textbf{廖融}\par\vspace{0.5em}
高卧白雲鄉，崖泉對閣凉。\par
守貧無屬道，多病數求方。\par
耕犢驚雷斃，寒蕪入圃荒。\par
如何帝未夢，吟苦頂鋪霜。\par

\section{句  其一}
\textbf{廖融}\par\vspace{0.5em}
雲穿搗藥屋，雪壓釣魚船。\par

\section{句  其二}
\textbf{廖融}\par\vspace{0.5em}
松風吹髮亂，巖溜濺碁寒。\par

\section{句  其三}
\textbf{廖融}\par\vspace{0.5em}
圭竇先知曉，盆池別見天。\par

\section{存目}
\textbf{廖融}\par\vspace{0.5em}
詩題：酬皇甫冉西陵見寄詩：西陵潮已漏，島嶼沒中流。\par
越客依風水，相看南渡頭。\par
寒光生極浦，暮雪映滄州。\par
何事揚帆去，空驚海上鷗。\par

\section{宮詞}
\textbf{翁宏}\par\vspace{0.5em}
又是春殘也，如何出翠帷。\par
落花人獨立，微雨燕雙飛。\par
寓目魂將斷，經年夢亦非。\par
那堪嚮秋夕，蕭颯暮蟾輝。\par

\section{秋風曲}
\textbf{翁宏}\par\vspace{0.5em}
又是秋殘也，無聊意若何。\par
客程江外遠，歸思夜深多。\par
峴首飛黄葉，湘湄走白波。\par
仍聞漢都護，今歲合休戈。\par

\section{贈衡山廖融南遊}
\textbf{翁宏}\par\vspace{0.5em}
病卧瘴雲間，莓苔漬竹關。\par
孤吟牛渚月，老憶洞庭山。\par
壯志潛消盡，淳風竟未還。\par
今朝忽相遇，執手一開顔。\par

\section{句  其一}
\textbf{翁宏}\par\vspace{0.5em}
風高弓力大，霜重角聲干。\par

\section{句  其二}
\textbf{翁宏}\par\vspace{0.5em}
客帆來異域，別島落蟠桃。\par

\section{句  其三}
\textbf{翁宏}\par\vspace{0.5em}
寒清萬國土，冷鬭四維根。\par

\section{句  其四}
\textbf{翁宏}\par\vspace{0.5em}
漏光殘井甃，缺影背山椒。\par

\section{句  其五}
\textbf{翁宏}\par\vspace{0.5em}
萬木殘秋裏，孤舟半夜猿。\par

\section{句  其六}
\textbf{翁宏}\par\vspace{0.5em}
因尋買珠客，誤入射猿家。\par

\section{句  其七}
\textbf{翁宏}\par\vspace{0.5em}
何處殘春夜，和花落古宮。\par

\section{句  其八}
\textbf{翁宏}\par\vspace{0.5em}
孤舟半夜雨，上國十年心。\par

\section{懷翁宏}
\textbf{王元}\par\vspace{0.5em}
獨夜思君切，無人知此情。\par
滄洲歸未得，華髮別來生。\par
孤館木初落，高空月正明。\par
遠書多隔歲，猶念沒前程。\par

\section{登祝融峰}
\textbf{王元}\par\vspace{0.5em}
草疊到孤頂，身齊高鳥翔。\par
勢疑撞翼軫，翠欲滴瀟湘。\par
雲濕幽崖滑，風梳古木香。\par
晴空聊縱目，杳杳極窮荒。\par

\section{聽琴}
\textbf{王元}\par\vspace{0.5em}
拂塵開素匣，有客獨傷時。\par
古調俗不樂，正聲君自知。\par
寒泉出澗澀，老檜倚風悲。\par
縱有來聽者，誰堪繼子期。\par

\section{題鄧真人遺址}
\textbf{王元}\par\vspace{0.5em}
三千功滿輕升去，留得山前舊隠基。\par
但見白雲長掩映，不知浮世幾興衰。\par
松梢風觸霓旌動，棕葉霜沾鶴翅垂。\par
近代無人尋異事，野泉噴月瀉秋池。\par

\section{悼李韶}
\textbf{王元}\par\vspace{0.5em}
韶也命何奇，生前與世違。\par
貧棲古梵剎，終著舊麻衣。\par
雅句僧抄遍，孤墳客吊稀。\par
故園今孰在，應見夢中歸。\par

\section{句  其一}
\textbf{王元}\par\vspace{0.5em}
伴行唯瘦鶴，尋寺入深雲。\par

\section{句  其二}
\textbf{王元}\par\vspace{0.5em}
江城賣藥常將鶴，古寺看碑不下馿。\par

\section{題全義分水嶺}
\textbf{伍彬}\par\vspace{0.5em}
前賢功及物，禹後杳難儔。\par
不及古今色，平分南北流。\par
寒衝山影岸，清繞荻花洲。\par
盡是朝宗去，潺湲早晚休。\par

\section{句  其一}
\textbf{伍彬}\par\vspace{0.5em}
稚子出看莎徑沒，漁翁來報竹橋流。\par

\section{句  其二}
\textbf{伍彬}\par\vspace{0.5em}
踪跡未辭鴛鷺客，夢魂先到鷓鴣村。\par

\section{贈藴上人}
\textbf{王正己}\par\vspace{0.5em}
佛法詩名誰更繼，未聞隨分謁侯王。\par
洗盂秋澗日華動，搗藥夜堂雲氣香。\par
苔蘚亂青封疊石，杉松濃影過空牆。\par
若非火嶽□□地，那得吾師住久長。\par

\section{贈廖融}
\textbf{王正己}\par\vspace{0.5em}
湖南雖不就卑官，高卧深雲道自安。\par
病起坐當秋閣迥，酒醒吟對夜濤寒。\par
爐中藥熟分僧服，榻上琴閑借我彈。\par
幸遇清朝有良鑒，退身爭忍似方干。\par

\section{偈}
\textbf{豐禪師}\par\vspace{0.5em}
駿馬機前異，遊人肘後懸。\par
既參雲外客，試爲老僧看。\par

\section{句}
\textbf{豐禪師}\par\vspace{0.5em}
雪嶺梅花綻，雲洞老僧驚。\par

\section{壽陽山}
\textbf{石仲元}\par\vspace{0.5em}
平原翠削萬瓊瑰，頓轡塵沙眼暫開。\par
文網牽人寧底急，未妨特特看山來。\par

\section{句  其一}
\textbf{石仲元}\par\vspace{0.5em}
地連錦野東西去，水接朱川次第來。\par

\section{句  其二}
\textbf{石仲元}\par\vspace{0.5em}
石壓木斜出，岸懸花倒生。\par

\section{吟贈夢英大師}
\textbf{王著}\par\vspace{0.5em}
到處聞人乞篆蹤，學來年久有深功。\par
墨池闊類湘江水，筆塚高齊太華峰。\par
金錫罷飛新解虎，鐵盂收掌舊降龍。\par
知師吟戀烟村景，不肯迴頭望九重。\par

\section{禁林讌會之什}
\textbf{梁周翰}\par\vspace{0.5em}
寶書驚絕耀天章，飛白親題賜玉堂。\par
瑞彩上騰流素月，朗河下注映丹牆。\par
鶴盤吳嶼雙翎健，鵲顧雕陵巨翼長。\par
遊霧半收懸組練，輕雲斜拂駐鸞皇。\par
墨池併獲三奇寶，翠琰俱生五色光。\par
陪讌禁林知有幸，叩頭遥祝萬年觴。\par

\section{題義門胡氏華林書院}
\textbf{梁周翰}\par\vspace{0.5em}
知有英儒伯始孫，聚書講學賁邱園。\par
征搜合副來賢詔，旌表宜題遍德門。\par
舉族爲人俱孝弟，深幽何草不蘭蓀。\par
小冠子夏來相示，詩版因憑寄竹軒。\par

\section{句  其一}
\textbf{梁周翰}\par\vspace{0.5em}
誰似金華楊學士，十聯詩在御屏間。\par

\section{句  其二}
\textbf{梁周翰}\par\vspace{0.5em}
百花將盡牡丹坼，十雨初晴太液春。\par

\section{句  其三}
\textbf{梁周翰}\par\vspace{0.5em}
宿雨一番蔬甲拆，春山幾處茗旗香。\par
ff。\par

\section{寄懷英公大師清交}
\textbf{馬去非}\par\vspace{0.5em}
雲隠秋鴻水隠魚，相思難得惠休書。\par
遥聞養性棲蓮岳，不肯携筇入帝都。\par
金殿聖緣應未斷，玉堂知己漸凋疏。\par
何人曾得陪高論，頭戴神羊馬大夫。\par

\section{宿山寺}
\textbf{孟貫}\par\vspace{0.5em}
溪山盡日行，方聽遠鐘聲。\par
入院逢僧定，登樓見月生。\par
露垂羣木潤，泉落一巖清。\par
此景關吾事，通宵寐不成。\par

\section{贈棲隠洞譚先生}
\textbf{孟貫}\par\vspace{0.5em}
先生雙鬢華，深谷卧雲霞。\par
不伐有巢樹，多移無主花。\par
石泉春釀酒，松火夜煎茶。\par
因問山中事，如君有幾家。\par

\section{歸雁}
\textbf{孟貫}\par\vspace{0.5em}
春至衡陽雁，思歸塞路長。\par
汀洲齊奮翼，霄漢共成行。\par
雪盡翻風暖，寒收度月凉。\par
直應到秋日，依舊返瀟湘。\par

\section{春江送人}
\textbf{孟貫}\par\vspace{0.5em}
春江多去情，相去枕長汀。\par
數雁別湓浦，片帆離洞庭。\par
雨餘沙草綠，雲散岸峰青。\par
誰共觀明月，漁歌夜好聽。\par

\section{過秦嶺}
\textbf{孟貫}\par\vspace{0.5em}
古今傳此嶺，高下勢崢嶸。\par
安得青山路，化爲平地行。\par
蒼苔留虎跡，碧樹障溪聲。\par
欲過一回首，踟躕無限情。\par

\section{送吳夢誾歸閩}
\textbf{孟貫}\par\vspace{0.5em}
甌閩在天末，此去整行衣。\par
久客逢春盡，思家冒暑歸。\par
海雲添晚景，山瘴滅晴暉。\par
相憶吟偏苦，不堪書信稀。\par

\section{山中夏日}
\textbf{孟貫}\par\vspace{0.5em}
深山宜避暑，門戶映嵐光。\par
夏木蔭溪路，晝雲埋石床。\par
心源澄道靜，衣葛蘸泉凉。\par
算得紅塵裏，誰知此興長。\par

\section{宿故人江居}
\textbf{孟貫}\par\vspace{0.5em}
渡口樹冥冥，南山漸隠青。\par
漁舟歸舊浦，鷗鳥宿前汀。\par
靜榻懸燈坐，閒門對浪扃。\par
相思頻到此，幾番醉還醒。\par

\section{寄伍喬}
\textbf{孟貫}\par\vspace{0.5em}
蹉跎春又晚，天末信來遲。\par
長憶分携日，正當摇落時。\par
獨遊饒旅恨，多事失歸期。\par
君看前溪樹，山禽巢幾枝。\par

\section{山中訪人不遇}
\textbf{孟貫}\par\vspace{0.5em}
負琴兼杖藜，特地過巖西。\par
已見竹軒閉，又聞山鳥啼。\par
長松寒倚谷，細草暗連溪。\par
久立無人事，烟霞歸路迷。\par

\section{贈隠者}
\textbf{孟貫}\par\vspace{0.5em}
世路爭名利，深山獨結茅。\par
安情自得所，非道豈相交。\par
百尺松當戶，千年鶴在巢。\par
知君於此景，未欲等閒拋。\par

\end{document}